% ============ Criterio: datos ============
\section{Datos utilizados: suficiencia y corrección}

VERIFICADO EMPÍRICAMENTE: el consolidado versionado (XLSX y CSV en datos/tarea\_join/) contiene solo 50.891 filas y 3 convocatorias (2017: 13.001, 2019: 16.796, 2021: 21.094), frente a las 77.237 filas y 6 convocatorias (2013-2021) declaradas en el README y el catálogo. Las evidencias versionadas (p. ej. matrices de transición 2013-2014) sí cubren las 6 convocatorias, lo que indica que el artefacto versionado es un snapshot de la fase inicial de 3 convocatorias nunca actualizado. Consecuencia verificada: ejecutar el pipeline en un clon fresco genera solo 2 periodos de transición (2017-2019, 2019-2021) en vez de 5, y los resultados no reproducen el README. El CSV paralelo además presenta mojibake (UTF-8/Latin-1). El dataset de producción (\textasciitilde{}1.2 GB) está gitignored y no es reproducible desde el clon. En cuanto a la corrección de las fuentes: el padrón es la fuente oficial del sistema de reconocimiento y el dataset de producción su complemento natural; la elección de datos abiertos (datos.gov.co/Socrata) es apropiada y verificable.


\subsection{Fuentes del repositorio}

\begin{itemize}
  \item \textbf{Padrón de investigadores} (Socrata \texttt{bqtm-4y2h}):
        convocatorias 2013--2021, 30 variables, \textbf{77\,237 registros
        declarados} y 30\,086 investigadores únicos.
  \item \textbf{Producción de grupos} (Socrata \texttt{33dq-ab5a}):
        3\,166\,629 filas producto-autor; dataset grande, no versionado en git.
  \item \textbf{Catálogo}: \texttt{datos/catalogo.yaml} con metadatos,
        diccionario de datos y llaves de cruce (\texttt{id\_persona\_pr} ↔
        \texttt{id\_persona\_pd}).
\end{itemize}

\begin{tcolorbox}[nota]
Los datos no son solo los archivos: también se evaluó \emph{cómo se
versionaron} en el repositorio. El detalle de la verificación (conteos reales,
codificación, hallazgos) está en este documento y en el maestro.
\end{tcolorbox}
